\documentclass[12pt]{article}
\usepackage[utf8]{inputenc}
\usepackage[T1]{fontenc}
\usepackage{hyperref}
\title{Pyen Language Resources for K-12 Curriculum Development}
\author{Compiled by Claude Sonnet 5.5 (claude-sonnet-5-5) for Andr\'as Kornai}
\date{}
\begin{document}
\maketitle

\section{Introduction}
Pyen (ISO 639-3 pyy; Glottocode pyen1239) is a Loloish (Sino-Tibetan) language of Shan State, Myanmar, with an estimated 600 to 700 speakers in two villages near Mong Yang. It is closely related to Bisu, Laomian and Laopin. Only a few academic resources turned up in this search, and no teaching materials.

\section{Primary Educational Resources}
None found.

\section{Linguistic and Reference Materials}
\begin{itemize}
\item Wikipedia article on Pyen (cites Kirk Person's 2007 phonological sketch): \url{https://en.wikipedia.org/wiki/Pyen_language}
\item Christina Scotte Horney, ``Tonal Variation in Pyen'' (2019): \url{https://evols.library.manoa.hawaii.edu/items/34a2dc8f-aca7-4e8a-af58-720403ea04c8}
\item Journal of the Southeast Asian Linguistics Society issue (Feb 2019), listing the tone study: \url{https://doaj.org/article/eeed001180384805b2315a3ebc691373}
\end{itemize}

\section{Religious and Community Resources}
None found.

\section{Digital and Audio Resources}
None found.

\section{Curriculum Development Context}
The community is small, and the search found no orthography, school programme or written materials. Any curriculum would need documentation work with the community first, building on the phonological and tonal descriptions above.
\end{document}